\chapter{Baseband File System}
\label{chap: baseband}

This chapter is here to give an overview of (and hopefully in more polished versions) give some technical details
into how we store and manipulate baseband data. Not super straightforward, and took me a while to understand.

\begin{itemize}
\item Baseband files are organized by UTC time. 
First the split is by 5 numbers (e.g. 17219), and then split again into raw files of length of about 45 seconds (in the MARS 1 case). 
They have a length of 10 numbers, corresponding to the UTC time at which the measurement of their data \textit{starts}. 
These raw files then contain the FTed data for all our channels, in complex numbers.

\item also note that there are two ways for us to store data, corresponding to different number of bits: 1 bit complex or 4 bit complex 
(corresponding to 2 bits total, or 8 bits total respectively).

\item Since it's an insane amount of data, it is compressed. The way in which it is compressed is in an 8-digit binary, corresponding to either an integer or a complex number. 
unpacking into float complex 64 takes a bunch of time right (explain better too lazy to clean up right away) ('uint8' vs 'complex64')

\item It may also LOSE DATA, so some spectra are missing. This isn't great, but there are functions to help remedy this. 
Make continuous and such are very valuable. Give brief overview, and perhaps link the function page itself.

\end{itemize}